\documentclass[10pt,a4paper]{amsart}
\usepackage{amsmath,amssymb,amsthm,graphicx}
% Retained preamble for the Stacks Project chapter snapshot in this directory.
% The chapter's own preamble.tex is a copy of the shared Stacks Project preamble; the block
% below is sliced MECHANICALLY, line for line, from it. Nothing was added, removed or
% redefined; the wrapper class, the xr-hyper/hyperref wiring and the externaldocument list
% are omitted because the assembler supplies the document class.
%
% lines taken from preamble.tex: 3, 5, 9, 11, 12, 156, 157, 158, ... (51 lines total)
%
% package coverage restored from the chapter's own preamble (2026-10-13): xy, xy-2cell

\usepackage[all]{xy}
\xyoption{2cell}
\UseAllTwocells

\newenvironment{reference}{\comment}{\endcomment}
\newenvironment{slogan}{\comment}{\endcomment}
\newenvironment{history}{\comment}{\endcomment}
\theoremstyle{plain}
\newtheorem{theorem}[subsection]{Theorem}
\newtheorem{proposition}[subsection]{Proposition}
\newtheorem{lemma}[subsection]{Lemma}
\theoremstyle{definition}
\newtheorem{definition}[subsection]{Definition}
\newtheorem{example}[subsection]{Example}
\newtheorem{exercise}[subsection]{Exercise}
\newtheorem{situation}[subsection]{Situation}
\theoremstyle{remark}
\newtheorem{remark}[subsection]{Remark}
\newtheorem{remarks}[subsection]{Remarks}
\numberwithin{equation}{subsection}
\def\lim{\mathop{\mathrm{lim}}\nolimits}
\def\colim{\mathop{\mathrm{colim}}\nolimits}
\def\Spec{\mathop{\mathrm{Spec}}}
\def\Hom{\mathop{\mathrm{Hom}}\nolimits}
\def\Ext{\mathop{\mathrm{Ext}}\nolimits}
\def\SheafHom{\mathop{\mathcal{H}\!\mathit{om}}\nolimits}
\def\SheafExt{\mathop{\mathcal{E}\!\mathit{xt}}\nolimits}
\def\Sch{\mathit{Sch}}
\def\Mor{\mathop{\mathrm{Mor}}\nolimits}
\def\Ob{\mathop{\mathrm{Ob}}\nolimits}
\def\Sh{\mathop{\mathit{Sh}}\nolimits}
\def\NL{\mathop{N\!L}\nolimits}
\def\CH{\mathop{\mathrm{CH}}\nolimits}
\def\proetale{{pro\text{-}\acute{e}tale}}
\def\etale{{\acute{e}tale}}
\def\QCoh{\mathit{QCoh}}
\def\Ker{\mathop{\mathrm{Ker}}}
\def\Im{\mathop{\mathrm{Im}}}
\def\Coker{\mathop{\mathrm{Coker}}}
\def\Coim{\mathop{\mathrm{Coim}}}
\DeclareMathSymbol{\boxtimes}{\mathbin}{AMSa}{"02}
\def\QCohstack{\mathcal{QC}\!\mathit{oh}}
\def\Cohstack{\mathcal{C}\!\mathit{oh}}
\def\Spacesstack{\mathcal{S}\!\mathit{paces}}
\def\Quotfunctor{\mathrm{Quot}}
\def\Hilbfunctor{\mathrm{Hilb}}
\def\Curvesstack{\mathcal{C}\!\mathit{urves}}
\def\Polarizedstack{\mathcal{P}\!\mathit{olarized}}
\def\Complexesstack{\mathcal{C}\!\mathit{omplexes}}
\def\Pic{\mathop{\mathrm{Pic}}\nolimits}
\def\Picardstack{\mathcal{P}\!\mathit{ic}}
\def\Picardfunctor{\mathrm{Pic}}
\def\Deformationcategory{\mathcal{D}\!\mathit{ef}}
\usepackage{cleveref}
\crefname{originallemma}{original lemma}{original lemmas}
\Crefname{originallemma}{Original lemma}{Original lemmas}
\makeatletter
\newlabel{categories-lemma-finite-limits-exist}{{002O}{source}{}{sourcepointer-1}{}}
\newlabel{categories-lemma-finite-limits-exist@cref}{{[originallemma][0][]002O}{[1][0][]source}{}{}{}}
\makeatother
\crefname{originallemma}{original lemma}{original lemmas}
\Crefname{originallemma}{Original lemma}{Original lemmas}
\makeatletter
\newlabel{definition-covering-SR}{{01G3}{source}{}{sourcepointer-2}{}}
\newlabel{definition-covering-SR@cref}{{[originallemma][0][]01G3}{[1][0][]source}{}{}{}}
\makeatother
\crefname{originallemma}{original lemma}{original lemmas}
\Crefname{originallemma}{Original lemma}{Original lemmas}
\makeatletter
\newlabel{lemma-basis-hypercovering}{{01H6}{source}{}{sourcepointer-3}{}}
\newlabel{lemma-basis-hypercovering@cref}{{[originallemma][0][]01H6}{[1][0][]source}{}{}{}}
\makeatother
\crefname{originallemma}{original lemma}{original lemmas}
\Crefname{originallemma}{Original lemma}{Original lemmas}
\makeatletter
\newlabel{lemma-split}{{0DAU}{source}{}{sourcepointer-4}{}}
\newlabel{lemma-split@cref}{{[originallemma][0][]0DAU}{[1][0][]source}{}{}{}}
\makeatother
\title{Existence of hypercoverings with members in a subset B}
\author{The Stacks Project Authors}
\date{stacks-project:hypercovering.tex:local-snapshot-20261012; GNU FDL 1.2 or later}
\begin{document}
\maketitle
\noindent\textit{Editorial note.} Statement span is lines 2796-2812 of sources/stacks-hypercovering/pinned-original-source.tex; proof span is lines 2814-2954, the complete author proof (139 body lines).\par\medskip
\noindent\textit{Editorial note.} Navigation stubs: 4 cited-but-not-included labels are supplied as non-load-bearing pointers with their pinned tags where the snapshot tag index has them.\par\medskip
\noindent\textit{Editorial note.} extra\_packages is empty on purpose: the retained preamble already loads xy with options and a second declaration clashes in the assembler wrapper.\par\medskip
\noindent\textit{Editorial note.} Text is retained verbatim, including the author's own typos and placeholders.\par\medskip
\noindent\textit{Editorial note.} Queue index 0 of queue\_hypercovering.json (pb=5278, stmt 2796-2812, proof 2814-2954).\par\medskip
\noindent\textit{Editorial note.} stacks-hypercovering wave-1 (hc01): fresh adjudication accept (2026-10-14); pool per standing fallback order (hypercovering -> pic -> spaces-more-cohomology) after md01 closed the moduli pool; hypercovering.tex 53 labels/34 with proof judged 26 accept / 8 reject / 14 pb=0 env-dup rows excluded; spans cite-free and reference/slogan/history-free so no patchenv needed. Named omission: Author omits the 'straightforward but cumbersome' verification that the maps K(phi) compose correctly - a named mechanical check; the construction and the covering verification are in-span.\par\medskip
\noindent\textit{Allowed standard prerequisites (editorial).}\begin{itemize}
\item Explicit prerequisite boundary. The cited inputs are prerequisites; the argument in the retained span is not replaced by any cited result.
\end{itemize}\medskip
\ifdefined\hypertarget\hypertarget{sourcepointer-1}{}\fi
\noindent\textit{Source reference.} Original lemma 002O: sources/stacks-hypercovering/pinned-original-source.tex, source lines 2954--2954. Cited by the retained statement or proof but not part of the collected spans; supplied only as a navigation pointer. The argument does not depend on its content.\par\medskip
\ifdefined\hypertarget\hypertarget{sourcepointer-2}{}\fi
\noindent\textit{Source reference.} Original lemma 01G3: sources/stacks-hypercovering/pinned-original-source.tex, source lines 510--521. Cited by the retained statement or proof but not part of the collected spans; supplied only as a navigation pointer. The argument does not depend on its content.\par\medskip
\ifdefined\hypertarget\hypertarget{sourcepointer-3}{}\fi
\noindent\textit{Source reference.} Original lemma 01H6: sources/stacks-hypercovering/pinned-original-source.tex, source lines 2655--2660. Cited by the retained statement or proof but not part of the collected spans; supplied only as a navigation pointer. The argument does not depend on its content.\par\medskip
\ifdefined\hypertarget\hypertarget{sourcepointer-4}{}\fi
\noindent\textit{Source reference.} Original lemma 0DAU: sources/stacks-hypercovering/pinned-original-source.tex, source lines 2767--2777. Cited by the retained statement or proof but not part of the collected spans; supplied only as a navigation pointer. The argument does not depend on its content.\par\medskip
\begin{lemma}
\label{lemma-hypercovering-object}
Let $\mathcal{C}$ be a site with fibre products.
Let $\mathcal{B} \subset \Ob(\mathcal{C})$ be a subset.
Assume
\begin{enumerate}
\item any object $U$ of $\mathcal{C}$ has a covering
$\{U_j \to U\}_{j \in J}$ with $U_j \in \mathcal{B}$, and
\item if $\{U_j \to U\}_{j \in J}$ is a covering
with $U_j \in \mathcal{B}$ and $\{U' \to U\}$ is a morphism with
$U' \in \mathcal{B}$, then $\{U_j \to U\}_{j \in J} \amalg \{U' \to U\}$
is a covering.
\end{enumerate}
Then for any $X$ in $\mathcal{C}$ there is a hypercovering $K$
of $X$ such that $K_n = \{U_{n, i}\}_{i \in I_n}$
with $U_{n, i} \in \mathcal{B}$ for all $i \in I_n$.
\end{lemma}

\begin{proof}
A warmup for this proof is the proof of
Lemma \ref{lemma-basis-hypercovering} and
we encourage the reader to read that proof first.

\medskip\noindent
First we replace $\mathcal{C}$ by the site $\mathcal{C}/X$.
After doing so we may assume that $X$ is the final object
of $\mathcal{C}$ and that $\mathcal{C}$ has all finite limits
(Categories, Lemma \ref{categories-lemma-finite-limits-exist}).

\medskip\noindent
Let $n \geq 0$. Let us say that an
{\it $n$-truncated $\mathcal{B}$-hypercovering of $X$}
is given by an $n$-truncated simplicial object $K$
of $\text{SR}(\mathcal{C})$
such that for $i \in I_a$, $0 \leq a \leq n$
we have $U_{a, i} \in \mathcal{B}$ and such that
$K_0$ is a covering of $X$ and
$K_{a + 1} \to (\text{cosk}_a \text{sk}_a K)_{a + 1}$
for $a = 0, \ldots, n - 1$
is a covering as in Definition \ref{definition-covering-SR}.

\medskip\noindent
Since $X$ has a covering $\{U_{0, i} \to X\}_{i \in I_0}$
with $U_i \in \mathcal{B}$ by assumption, we get a $0$-truncated
$\mathcal{B}$-hypercovering of $X$. Observe that any $0$-truncated
$\mathcal{B}$-hypercovering of $X$ is split, see
Lemma \ref{lemma-split}.

\medskip\noindent
The lemma follows if we can prove for $n \geq 0$ that given a
split $n$-truncated $\mathcal{B}$-hypercovering $K$ of $X$
we can extend it to a
split $(n + 1)$-truncated $\mathcal{B}$-hypercovering of $X$.

\medskip\noindent
Construction of the extension. Consider the $(n + 1)$-truncated simplicial
object $K' = \text{sk}_{n + 1}(\text{cosk}_n K)$ of $\text{SR}(\mathcal{C})$.
Write
$$
K'_{n + 1} = \{U'_{n + 1, i}\}_{i \in I'_{n + 1}}
$$
Since $K = \text{sk}_n K'$ we have $K_a = K'_a$ for $0 \leq a \leq n$.
For every $i' \in I'_{n + 1}$ we choose a covering
\begin{equation}
\label{equation-choose-covering-B}
\{g_{n + 1, j} : U_{n + 1, j} \to U'_{n + 1, i'}\}_{j \in J_{i'}}
\end{equation}
with $U_{n + 1, j} \in \mathcal{B}$ for $j \in J_{i'}$.
This is possible by our assumption on $\mathcal{B}$ in the lemma.
For $0 \leq m \leq n$ denote $N_m \subset I_m$ the subset of
nondegenerate indices. We set
$$
I_{n + 1} =
\coprod\nolimits_{\varphi : [n + 1] \to [m]\text{ surjective, }0\leq m \leq n}
N_m \amalg
\coprod\nolimits_{i' \in I'_{n + 1}} J_{i'}
$$
For $j \in I_{n + 1}$ we set
$$
U_{n + 1, j} =
\left\{
\begin{matrix}
U_{m, i} & \text{if} &
j = (\varphi, i) & \text{where} & \varphi : [n + 1] \to [m], i \in N_m \\
U_{n + 1, j} & \text{if} &
j \in J_{i'} & \text{where} & i' \in I'_{n + 1}
\end{matrix}
\right.
$$
with obvious notation. We set $K_{n + 1} = \{U_{n + 1, j}\}_{j \in I_{n + 1}}$.
By construction $U_{n + 1, j}$ is an element
of $\mathcal{B}$ for all $j \in I_{n + 1}$. Let us define compatible
maps
$$
I_{n + 1} \to I'_{n + 1}
\quad\text{and}\quad
K_{n + 1} \to K'_{n + 1}
$$
Namely, the first map is given by
$(\varphi, i) \mapsto \alpha'(\varphi)(i)$ and
$(j \in J_{i'}) \mapsto i'$.
For the second map we use the morphisms
$$
f'_{\varphi, i} : U_{m, i} \to U'_{n + 1, \alpha'(\varphi)(i)}
\quad\text{and}\quad
g_{n + 1, j} : U_{n + 1, j} \to U'_{n + 1, i'}
$$
We claim the morphism
$$
K_{n + 1} \to K'_{n + 1} =
(\text{cosk}_n \text{sk}_n K')_{n + 1} =
(\text{cosk}_n K)_{n + 1}
$$
is a covering as in Definition \ref{definition-covering-SR}.
Namely, if $i' \in I'_{n + 1}$, then either $i'$ is nondegenerate
and the inverse image of $i'$ in $I_{n + 1}$ is equal to $J_{i'}$
and we get a covering of $U'_{n + 1, i'}$ by our choice
(\ref{equation-choose-covering-B}), or $i'$ is degenerate and
the inverse image of $i'$ in $I_{n + 1}$ is
$J_{i'} \amalg \{(\varphi, i)\}$ for a unique pair $(\varphi, i)$
and we get a covering by our choice (\ref{equation-choose-covering-B})
and assumption (2) of the lemma.

\medskip\noindent
To finish the proof we have to define the morphisms
$K(\varphi) : K_{n + 1} \to K_m$ corresponding to morphisms
$\varphi : [m] \to [n + 1]$, $0 \leq m \leq n$ and the morphisms
$K(\varphi) : K_m \to K_{n + 1}$ corresponding to morphisms
$\varphi : [n + 1] \to [m]$, $0 \leq m \leq n$
satisfying suitable composition relations.
For the first kind we use the composition
$$
K_{n + 1} \to K'_{n + 1} \xrightarrow{K'(\varphi)} K'_m = K_m
$$
to define $K(\varphi) : K_{n + 1} \to K_m$.
For the second kind, suppose given $\varphi : [n + 1] \to [m]$,
$0 \leq m \leq n$. We define the corresponding morphism
$K(\varphi) : K_m \to K_{n + 1}$ as follows:
\begin{enumerate}
\item for $i \in I_m$ there is a unique surjective map
$\psi : [m] \to [m_0]$ and a unique $i_0 \in I_{m_0}$ nondegenerate
such that $\alpha(\psi)(i_0) = i$\footnote{For example, if $i$ is
nondegenerate, then $m = m_0$ and $\psi = \text{id}_{[m]}$.},
\item we set $\varphi_0 = \psi_0 \circ \varphi : [n + 1] \to [m_0]$
and we map
$i \in I_m$ to $(\varphi_0, i_0) \in I_{n + 1}$, in other words,
$\alpha(\varphi)(i) = (\varphi_0, i_0)$, and
\item the morphism
$f_{\varphi, i} : U_{m, i} \to U_{n + 1, \alpha(\varphi)(i)} = U_{m_0, i_0}$
is the inverse of the isomorphism $f_{\psi, i_0} : U_{m_0, i_0} \to U_{m, i}$
(see Lemma \ref{lemma-split}).
\end{enumerate}
We omit the straightforward but cumbersome verification that this defines
a split $(n + 1)$-truncated $\mathcal{B}$-hypercovering of $X$
extending the given $n$-truncated one. In fact, everything is clear
from the above, except for the verification that the morphisms
$K(\varphi)$ compose correctly for all $\varphi : [a] \to [b]$
with $0 \leq a, b \leq n + 1$.
\end{proof}

\end{document}
